% =====================================================================
%  PART III — COMBINED APPENDIX
% =====================================================================
% AUTO-GENERATED by build/gen_tables — do not edit by hand
\newcommand{\ambforecast}{%
\begin{tabular}{@{}p{0.235\textwidth}rrrrrrrrrr@{}}
\hrow \wh{Rs crore unless stated} & \wh{FY27} & \wh{FY28} & \wh{FY29} & \wh{FY30} & \wh{FY31} & \wh{FY32} & \wh{FY33} & \wh{FY34} & \wh{FY35} & \wh{FY36}\\
Capacity (Mn.T) & 112 & 116 & 119 & 121 & 123 & 125 & 127 & 129 & 131 & 133\\
\arow Utilisation & 70.0\% & 72.5\% & 75.0\% & 77.5\% & 79.5\% & 81.0\% & 82.5\% & 83.5\% & 84.5\% & 85.0\%\\
Volume (Mn.T) & 78.4 & 84.1 & 89.2 & 93.8 & 97.8 & 101.2 & 104.8 & 107.7 & 110.7 & 113.0\\
\arow Realisation (Rs/t) & 5,758 & 5,988 & 6,180 & 6,328 & 6,455 & 6,558 & 6,650 & 6,729 & 6,810 & 6,892\\
Revenue & 45,141 & 50,359 & 55,153 & 59,340 & 63,116 & 66,398 & 69,671 & 72,486 & 75,385 & 77,913\\
\arow EBITDA & 7,710 & 9,061 & 10,357 & 11,513 & 12,599 & 13,559 & 14,498 & 15,366 & 16,227 & 16,985\\
EBITDA margin & 17.1\% & 18.0\% & 18.8\% & 19.4\% & 20.0\% & 20.4\% & 20.8\% & 21.2\% & 21.5\% & 21.8\%\\
\arow EBITDA per tonne (Rs) & 983 & 1,077 & 1,160 & 1,228 & 1,288 & 1,339 & 1,384 & 1,427 & 1,466 & 1,502\\
Depreciation & 3,550 & 4,026 & 4,424 & 4,750 & 5,010 & 5,209 & 5,392 & 5,561 & 5,716 & 5,859\\
\arow EBIT & 4,161 & 5,035 & 5,934 & 6,763 & 7,589 & 8,350 & 9,106 & 9,805 & 10,511 & 11,126\\
NOPAT & 3,114 & 3,768 & 4,440 & 5,061 & 5,679 & 6,248 & 6,814 & 7,337 & 7,866 & 8,326\\
\arow Capex & 9,500 & 9,000 & 8,500 & 8,000 & 7,500 & 7,500 & 7,500 & 7,500 & 7,500 & 7,500\\
$\Delta$ working capital & -14 & -18 & -16 & -14 & -13 & -11 & -11 & -10 & -10 & -9\\
\trow \textbf{FCFF} & -2,823 & -1,189 & 380 & 1,825 & 3,202 & 3,968 & 4,718 & 5,407 & 6,092 & 6,693\\
\arow Opening invested capital & 67,208 & 73,145 & 78,101 & 82,161 & 85,398 & 87,875 & 90,155 & 92,252 & 94,181 & 95,956\\
ROIC incl. goodwill & 4.6\% & 5.2\% & 5.7\% & 6.2\% & 6.7\% & 7.1\% & 7.6\% & 8.0\% & 8.4\% & 8.7\%\\
\arow ROIC operating only & 7.0\% & 7.5\% & 8.1\% & 8.6\% & 9.1\% & 9.6\% & 10.1\% & 10.6\% & 11.0\% & 11.4\%\\
EVA & -1,911 & -1,931 & -1,822 & -1,662 & -1,412 & -1,124 & -817 & -533 & -223 & 35\\
\end{tabular}}
\newcommand{\utclforecast}{%
\begin{tabular}{@{}p{0.30\textwidth}rrrrr@{}}
\hrow \wh{Rs crore unless stated} & \wh{FY27} & \wh{FY28} & \wh{FY29} & \wh{FY30} & \wh{FY31}\\
Capacity (Mn.T) & 186 & 195 & 205 & 213 & 220\\
\arow Utilisation & 82.0\% & 83.5\% & 84.5\% & 85.5\% & 86.0\%\\
Volume (Mn.T) & 152.5 & 162.8 & 173.2 & 182.1 & 189.2\\
\arow EBITDA per tonne (Rs) & 1,185 & 1,256 & 1,306 & 1,345 & 1,379\\
Realisation (Rs/t) & 6,072 & 6,254 & 6,379 & 6,475 & 6,556\\
\arow Standalone revenue & 92,609 & 101,832 & 110,503 & 117,916 & 124,035\\
Subsidiary revenue & 6,722 & 7,092 & 7,447 & 7,782 & 8,093\\
\arow Total revenue & 99,331 & 108,924 & 117,949 & 125,698 & 132,128\\
EBITDA & 19,749 & 22,234 & 24,514 & 26,489 & 28,175\\
\arow EBITDA margin & 19.9\% & 20.4\% & 20.8\% & 21.1\% & 21.3\%\\
Depreciation & 5,031 & 5,257 & 5,471 & 5,649 & 5,793\\
\arow EBIT & 14,718 & 16,977 & 19,043 & 20,840 & 22,381\\
NOPAT & 11,014 & 12,704 & 14,250 & 15,595 & 16,748\\
\arow Capex & 9,500 & 9,500 & 9,000 & 8,500 & 8,500\\
$\Delta$ working capital & 2,150 & -359 & -338 & -290 & -240\\
\trow \textbf{FCFF} & 4,394 & 8,820 & 11,059 & 13,034 & 14,282\\
\arow Opening invested capital & 93,761 & 100,381 & 104,265 & 107,456 & 110,017\\
ROIC & 11.7\% & 12.7\% & 13.7\% & 14.5\% & 15.2\%\\
\arow EVA & 298 & 1,232 & 2,334 & 3,314 & 4,175\\
\end{tabular}}
\newcommand{\ambgrid}{%
\begin{tabular}{@{}lrrrrr@{}}
\hrow \wh{WACC $\backslash$ g} & \wh{2\%} & \wh{3\%} & \wh{4\%} & \wh{5\%} & \wh{6\%}\\
\arow 9\% & 163 & 173 & 186 & 206 & 240\\
10\% & 125 & 129 & 136 & 144 & 158\\
\arow 11\% & 95 & 97 & 100 & 103 & 108\\
12\% & 72 & 72 & 73 & 74 & 75\\
\arow 13\% & 53 & 53 & 52 & 52 & 51\\
\end{tabular}}
\newcommand{\utclgrid}{%
\begin{tabular}{@{}lrrrrr@{}}
\hrow \wh{WACC $\backslash$ g} & \wh{4.0\%} & \wh{5.0\%} & \wh{6.0\%} & \wh{6.5\%} & \wh{7.0\%}\\
9.50\% & 5,681 & 6,348 & 7,395 & 8,181 & 9,281\\
\arow 10.50\% & 4,674 & 5,060 & 5,617 & 6,000 & 6,493\\
11.43\% & 3,981 & 4,221 & 4,549 & 4,763 & 5,026\\
\arow 12.50\% & 3,371 & 3,512 & 3,697 & 3,812 & 3,948\\
13.50\% & 2,925 & 3,010 & 3,118 & 3,184 & 3,260\\
\end{tabular}}


\partpage{Part III}{Combined Appendix}{Supporting data, workings and disclosures for
both companies. Everything in Parts I and II is reproducible from what follows and
from the two accompanying Excel models.}

\runhead{Part III \textbar\ Combined Appendix}

% ---------------------------------------------------------------- A
\banner{Appendix A}{Comparative chapter}{ONE EASY, ONE DIFFICULT}{Ambuja Rs 99.4}{UltraTech Rs 5,024}{Both SELL}{31 March 2026}

\begin{paracol}{2}
\section*{A\quad Comparative chapter: one easy, one difficult}
\addtocounter{section}{0}

The course pairs an easy-to-value company with a difficult one. UltraTech is the easy
name and Ambuja the difficult one, and the distinction is not about size or quality
but about how much had to be reconstructed before any return measure became
meaningful.

\subsection*{A.1\quad What made UltraTech easy}

A ten-year audited standalone history that ties both ways in every year. Disclosed
grey cement volume and capacity, which allow a capacity-times-utilisation forecast
rather than an assumed growth rate. A reported effective tax rate of 26.1\% against a
statutory 25.168\%, so no migration path is required. Explicit management guidance on
capex and on a Rs 300 per tonne cost saving. One dominant product. The modelling
problems were ordinary ones: choosing a horizon, handling circularity, and deciding
what the terminal return should be.

\subsection*{A.2\quad What made Ambuja difficult}

Five distinct problems, each of which had to be solved before the model could produce
a number worth reading.

\subsubsection{Two years of consolidated history} The annual report provides
FY2025-26 and a restated FY2024-25. Earlier consolidated years exist only in broker
research. A return trend asserted from two points is thin, and the ten-year explicit
forecast rests on a two-point base.

\subsubsection{A restated year and a fifteen-month year} FY2024-25 was restated for
the Sanghi and Penna amalgamation, and the transition year to March 2023 covered
fifteen months. Any multi-year growth rate taken off the reported series without
adjustment is wrong.

\subsubsection{A tax line that is a net credit} Reported tax was a credit of
Rs 2,338 crore, so reported profit exceeds pre-tax profit and the effective rate is
minus 71\%. Every return in the report had to be recomputed at the statutory rate,
and the credit decomposed into three drivers with different economics.

\subsubsection{A merger in progress} At the valuation date the structure was a parent
with listed subsidiaries; afterwards it is one entity with a 13.7\% larger share
count. Both bases are shown.

\subsubsection{A capital base that is one-third goodwill} Rs 22,979 crore of goodwill
and intangibles, which must sit inside invested capital for the return measure but
must be excluded from terminal reinvestment.

\subsection*{A.3\quad How the five consistencies were preserved in each}

In both models the same debt policy — constant $D/V$ with continuous rebalancing — is
applied in the unlever, the relever, the discounting of the tax shield and the
construction of the net borrowing line, and in both the identity $\rho = \text{WACC}
+ K_d t (D/V)$ ties exactly. In both, every liability included in the cost of capital
is subtracted in the bridge and none is subtracted twice. Both carry a PARTIAL on
asset consistency and a PARTIAL on multiplier consistency, for different reasons:
Ambuja's depreciation base does not match its depreciation charge, while UltraTech's
India Cements stake has an unreconciled treatment.

\switchcolumn
\ex{app_roic.png}{\textbf{Exhibit C1.} Both forecasts assume returns improve. Ambuja's
never crosses the cost of capital on the goodwill-inclusive measure; UltraTech's
starts above it. Source: authors' models.}

\ex{app_gw.png}{\textbf{Exhibit C3.} The structural difference between the two
companies in one number. Source: FY2025-26 annual reports.}

\ex{app_gap.png}{\textbf{Exhibit C4.} Source: authors' models; NSE closes 30 March
2026.}

\begin{extab}{The two companies side by side}{Ambuja consolidated; UltraTech
standalone core with subsidiaries added. Source: authors' models.}
\begin{tabular}{@{}p{0.44\linewidth}rr@{}}
\hrow \wh{} & \wh{Ambuja} & \wh{UltraTech}\\
Classification & Difficult & Easy\\
\arow Years of usable history & 2 & 10\\
Explicit forecast horizon & 10 yr & 5 yr\\
\arow Basis & Consol. & S/A core\\
Volume FY26 (Mn.T) & 74.0 & 140.8\\
\arow Utilisation FY26 & 67.9\% & 81.6\% \\
EBITDA per tonne (Rs) & 886 & 1,097\\
\arow ROIC FY26, incl. goodwill & 3.3\% & 10.8\% \\
Goodwill \% of invested capital & 34.2\% & 8.9\% \\
\arow WACC & 11.36\% & 11.43\% \\
Terminal growth & 6.99\% & 6.99\% \\
\arow Terminal ROIC & 11.9\% & 15.9\% \\
Terminal value \% of EV & 72.3\% & 78.7\% \\
\arow Implied terminal EV/EBITDA & 4.65x & 7.52x\\
Traded EV/EBITDA & 17.0x & 21.1x\\
\arow Price / book & 1.67x & 4.24x\\
Return on equity & 3.8\% & 10.4\% \\
\trow \textbf{Value per share (Rs)} & \textbf{99.4} & \textbf{5,024}\\
\trow \textbf{Market price (Rs)} & \textbf{400.90} & \textbf{10,745}\\
\trow \textbf{Downside} & \textbf{$-$75.2\%} & \textbf{$-$53.2\%}\\
\end{tabular}
\end{extab}

\switchcolumn*
\subsection*{A.4\quad Why two SELLs are consistent rather than contradictory}

The obvious objection is that Ambuja trades \emph{below} the peer median on EV/EBITDA
while UltraTech trades above it, so they cannot both be expensive. The steady-state
identity resolves it. In a steady state, enterprise value over invested capital equals
$(\text{ROIC}-g)/(\text{WACC}-g)$. Applying it to each company's own terminal return
gives a justified EV/IC of 2.02x for UltraTech and 1.13x for Ambuja. At traded prices
the actual figures are 3.31x and 1.66x, with invested capital measured including
goodwill for both. So UltraTech's price is about 1.6 times what its own terminal
return justifies, and Ambuja's about 1.5 times — and Ambuja's terminal return is
itself the more speculative of the two, because it requires a crossing the company has
never achieved.

\begin{keybox}\boxtitle{The claim we are making}
We are not claiming Ambuja is mispriced against its sector. We are claiming the sector
is mispriced against the cash it generates, and that Ambuja is the member with the
weakest returns. On today's return on capital, Ambuja's justified multiple of invested
capital is \emph{negative}: at a 3.3\% return and 7.0\% growth every rupee reinvested
destroys value, so no positive multiple of invested capital is warranted at all.
\end{keybox}

\subsection*{A.5\quad The known inconsistency between the two models}

UltraTech holds 74.99\% of the listed India Cements and is modelled from a standalone
core with the subsidiary block carried explicitly. Ambuja holds 50.05\% of the listed
ACC and is modelled consolidated. The reason is data availability rather than
principle: UltraTech publishes a ten-year standalone history and no comparable
consolidated one, while Ambuja's consolidated statements were available and its
standalone accounts exclude more than a third of the business. The irony is that
UltraTech's stake is the more controlling of the two.

The right resolution before submission is to consolidate both or to value both by sum
of the parts. Until that is done this is a disclosed weakness rather than a hidden one,
and its effect is quantified: UltraTech's omitted subsidiary block is under 10\% of
group EBITDA, against 55\% at Ambuja.

\switchcolumn
\ex{app_evic.png}{\textbf{Exhibit C2.} The steady-state identity $EV/IC =
(\text{ROIC}-g)/(\text{WACC}-g)$ evaluated on each company's own numbers. Ambuja's
justified multiple on today's return is negative because its return sits below its
growth rate. Source: authors' models.}

\ex{app_mult.png}{\textbf{Exhibit C5.} The multiplier-consistency gap for both names.
In each case the multiple our model implies is far below both the traded and the peer
multiple, and in each case we take a documented position rather than adopt the peer
figure. Source: authors' models; peer annual reports.}

\ex{app_audit.png}{\textbf{Exhibit C6.} The five consistencies plus the ten-point
audit checklist plus the seven Appendix 5A sanity checks, as recorded for each
company. Source: authors' assessment.}

\end{paracol}
\clearpage

% ---------------------------------------------------------------- B
\section*{B\quad Model audit log and corrections applied}

The models supplied for this report were audited independently: both were rebuilt from
their input assumptions in Python, and every output was compared cell for cell. The
rebuild reproduced Ambuja's published value to the paisa and UltraTech's to within
2.5 paise, and reproduced Ambuja's sensitivity grid exactly. That agreement is what
gave the discrepancies below their weight — where the rebuild and the workbook differ,
the workbook contains an error.

\vspace{2mm}
{\footnotesize
\begin{tabular}{@{}p{0.035\textwidth}p{0.17\textwidth}p{0.075\textwidth}p{0.35\textwidth}p{0.30\textwidth}@{}}
\hrow \wh{\#} & \wh{Where} & \wh{Severity} & \wh{Finding} & \wh{Action taken}\\
1 & UTCL Sensitivity grid & HIGH & The two-way grid discounts P\&L row 17 (depreciation) instead of row 24 (FCFF), builds its terminal value off row 14 (revenue) instead of row 21 (NOPAT), and omits the non-controlling interest from the bridge. It prints Rs 27,314--117,876 per share against a base case of Rs 5,148. & Rebuilt. Exhibit U8 is the corrected grid. Fixed in the corrected workbook.\\
\arow 2 & UTCL Source-Data NOPAT & HIGH & Historical NOPAT is computed as EBIT $\times(1-\text{Assumptions!}C4)$, but the tax rate lives in \texttt{B4}; \texttt{C4} is empty, so NOPAT equals EBIT and the published ten-year ROIC series is pre-tax. & Corrected to \texttt{B4}. Exhibit U4 uses the after-tax series (FY26: 10.8\%, not 14.4\%).\\
3 & Both, beta relever & HIGH & Both workbooks state Harris--Pringle constant rebalancing and discount the tax shield at $\rho$, but relever beta with the fixed-debt Hamada formula. Ambuja's formula additionally carries a sign error, computing $(1+t)$ where even Hamada needs $(1-t)$. & Relevered as $\beta_U(1+D/E)$ throughout, consistent with the stated debt policy. Effect: Ambuja Rs 98.85 $\rightarrow$ Rs 99.42; UltraTech Rs 5,148 $\rightarrow$ Rs 5,024.\\
\arow 4 & AMB Source-Data & MED & \texttt{C32} (EBIT) is \texttt{=C29-C7} where \texttt{C29} is a blank label row, so EBIT prints $-$3,570 instead of 2,989. \texttt{C34} (effective tax rate) references the same blank cell and prints 100\%. & Repointed to \texttt{C31} and \texttt{C33}. Display-only cells; no valuation effect.\\
5 & AMB Terminal & MED & \texttt{B21}, the implied exit EV/EBITDA, is \texttt{=B19/B20} where \texttt{B19} is empty, so it prints 0. & Repointed to \texttt{B17}. Prints 4.65x, matching the FCFF sheet.\\
\arow 6 & AMB Relative Valuation & MED & \texttt{C48} hardcodes a DCF value of Rs 114.87 that no longer matches the live model, and the DCF-implied EV/EBITDA in \texttt{C54} is computed from it. & Replaced with a live link to the FCFF sheet.\\
7 & Both, narrative text & MED & Substantial prose in both workbooks is stale relative to the live cells: Ambuja's Summary conclusion cites Rs 122 and Rs 151 against live values of Rs 98.85 and Rs 131.42, a WACC of 11.03\% against 11.37\%, and a sensitivity range of Rs 116--134 against a live Rs 95--108; its Dashboard cites a terminal growth of 5.0\% at 46.46\% of nominal GDP against a live 6.99\% at 65\%. UltraTech's Summary cites Rs 4,750 and a 5\% terminal growth; its Dashboard cites Rs 4,747.56. & All narrative in this report is written from the live cells. Workbook prose refreshed in the corrected versions.\\
\arow 8 & UTCL terminal fade switch & MED & The audit log records Golden Check 6 as ``FIXED $\ldots$ switched ON by default,'' but the switch ships at 0 (off) and the Dashboard correctly records the check as a FAIL. The two statements contradict each other. & Reported as shipped: fade off, check FAILs. The faded value of Rs 4,560 is disclosed in Part II.\\
9 & UTCL India Cements & MED & The bridge adds Rs 6,739.51 crore of investments. The balance sheet reports Rs 12,541 crore, of which India Cements is Rs 8,970 crore; the Entity Build sheet derives ``other investments'' as Rs 3,570.83 crore. None of the three tie. India Cements' operations are also already inside the forecast's subsidiary block. & Left as modelled and flagged. Adding the stake on top would double-count roughly Rs 197 per share. Recorded as an open asset-consistency item.\\
\arow 10 & AMB Audit Log, Entity Build, Dashboard & LOW & The audit-log block, the tie-out check and the seven Golden Checks are each duplicated five times on their sheets. & De-duplicated in the corrected workbook.\\
\end{tabular}\par}

\vspace{3mm}
\begin{minipage}[t]{0.48\textwidth}
\begin{extab}{Correction bridge: Ambuja}{Source: authors' analysis.}
\begin{tabular}{@{}p{0.62\linewidth}r@{}}
\hrow \wh{Step} & \wh{Rs/share}\\
As received (workbook) & 98.85\\
\arow Hamada relever with correct $(1-t)$ & 99.99\\
\trow \textbf{Harris--Pringle relever, $\beta_U(1+D/E)$} & \textbf{99.42}\\
Items 4, 5, 6, 10 (display and staleness) & no effect\\
\trow \textbf{Reported in this document} & \textbf{99.42}\\
\end{tabular}
\end{extab}
\end{minipage}\hfill
\begin{minipage}[t]{0.48\textwidth}
\begin{extab}{Correction bridge: UltraTech}{Source: authors' analysis.}
\begin{tabular}{@{}p{0.62\linewidth}r@{}}
\hrow \wh{Step} & \wh{Rs/share}\\
As received (workbook) & 5,148.18\\
\arow Items 1, 2 (grid and historical ROIC) & no effect on DCF\\
\trow \textbf{Harris--Pringle relever, $\beta_U(1+D/E)$} & \textbf{5,024.02}\\
Memo: with terminal ROIC fade applied & 4,560\\
\trow \textbf{Reported in this document} & \textbf{5,024.02}\\
\end{tabular}
\end{extab}
\end{minipage}

\vspace{2mm}
\begin{keybox}\boxtitle{The recommendation does not turn on any of this}
Every correction above is immaterial to the conclusion. The largest moves Ambuja by
0.6\% and UltraTech by 2.4\%, against downside of 75\% and 53\% respectively. They are
reported in full because a valuation that cannot survive its own audit is not worth
presenting, not because they change the answer.
\end{keybox}
\clearpage

% ---------------------------------------------------------------- C
\section*{C\quad The five consistencies, side by side}

{\footnotesize
\begin{tabular}{@{}p{0.028\textwidth}p{0.10\textwidth}p{0.055\textwidth}p{0.355\textwidth}p{0.055\textwidth}p{0.355\textwidth}@{}}
\hrow \wh{\#} & \wh{Consistency} & \wh{AMB} & \wh{Ambuja: evidence and gap} & \wh{UTCL} & \wh{UltraTech: evidence and gap}\\
1 & Assumptions & \PASS & Volume is capacity $\times$ utilisation, so growth is tied to physical assets. Capex is an explicit input anchored to the expansion plan. Terminal year forced to steady state with reinvestment $= g \times$ operating invested capital. & \PASS & Same construction. Growth requires reinvestment throughout; capex is explicit and anchored to guidance; terminal year is a forced steady state.\\
\arow 2 & Risk & \PASS & FCFF at WACC; CCF and APV at $\rho$; FCFE at $K_e$; EVA at WACC. Constant $D/V$ rebalancing applied in unlever, relever and shield. $\rho = \text{WACC}+K_d t(D/V)$ ties. & \PASS & Identical treatment. At a 6.1\% $D/V$ the policy choice actually matters here, and it is applied consistently.\\
3 & Asset & \PART & Other income excluded from EBITDA; cash non-operating and added back once; cash-corrected beta. \textbf{Gap:} depreciation charged at 8\% of operating fixed assets while the historical charge includes amortisation of a Rs 9,433 cr intangible balance held flat. & \PART & Same treatment of cash and beta. \textbf{Gap:} the India Cements stake is at book, the investments total does not tie to the bridge, and the stake's operations are already inside the subsidiary block.\\
\arow 4 & Liability & \PASS & Borrowings and leases in the weights and in the bridge. Deferred tax in neither the weights nor invested capital, but deducted in the bridge. NCI deducted (worth Rs 51/share). & \PASS & Loan funds and leases in both. Deferred tax of Rs 8,506 cr in neither, consistently on both sides. NCI of Rs 4,089 cr deducted.\\
5 & Multiplier & \PART & Implied terminal EV/EBITDA 4.65x against a peer median of 19.4x. Peer multiple solved back implies $g$ of 10.99\%, above nominal GDP of 10.76\%. Position taken; gap real. & \PART & Implied 7.52x against 19.05x. Peer multiple implies $g$ of 9.8\%, or 91\% of nominal GDP in perpetuity. Position taken; gap real.\\
\end{tabular}\par}

\vspace{4mm}
\section*{C.2\quad The ten-point audit checklist}

{\footnotesize
\begin{tabular}{@{}p{0.028\textwidth}p{0.40\textwidth}p{0.055\textwidth}p{0.055\textwidth}p{0.36\textwidth}@{}}
\hrow \wh{\#} & \wh{Check} & \wh{AMB} & \wh{UTCL} & \wh{Note}\\
1 & Five methods reconcile within 0.5\% & \PASS & \PASS & Both reconcile to six decimal places, and both were reproduced independently.\\
\arow 2 & WACC recomputes; market-value weights; circularity handled & \PASS & \PASS & Ambuja: immaterial at 0.87\% D/E. UltraTech: iterated per SSRN 413240, 1.0\% difference.\\
3 & Debt grows at $g$ where constant $D/V$ assumed & \PASS & \PASS & Debt set to $D/V \times$ EV at every date; net new borrowing line visible on the FCFE sheet.\\
\arow 4 & Cash and investments: one treatment throughout & \PASS & \PART & UltraTech's India Cements treatment is the open item.\\
5 & Every liability in WACC subtracted in the bridge & \PASS & \PASS & Verified line by line in both.\\
\arow 6 & Terminal: $g<$ WACC, built multiplicatively, reinvestment $=g/$RONIC & \PASS & \PART & $g$ built as $(1+6.5\%)(1+4.0\%)-1$ taken at 65\% in both. UltraTech's terminal ROIC does not fade.\\
7 & Tax: marginal statutory rate & \PASS & \PASS & 25.168\% under s.115BAA in both. Ambuja's reported rate is unusable; UltraTech's 26.1\% is in line.\\
\arow 8 & No DDT on post-FY2020 flows & \PASS & \PASS & None applied. DDT abolished with effect from FY2020-21.\\
9 & Every input traces to the source log & \PART & \PART & Open items: Ambuja's merger terms; UltraTech's capacity, utilisation and capex guidance, and its derived standalone capacity of 172.5 Mn.T.\\
\arow 10 & Sensitivity tables live-linked & \PASS & \FAIL & Ambuja's grid is correct and live. UltraTech's was found broken (item B1) and has been rebuilt.\\
\end{tabular}\par}
\clearpage

% ---------------------------------------------------------------- D & E
\section*{D\quad Cost of capital workings, both companies}

\begin{minipage}[t]{0.48\textwidth}
\begin{extab}{Shared inputs}{Source: RBI/FBIL Ratios and Rates release 27-Mar-2026;
Damodaran, NYU Stern, January 2026.}
\begin{tabular}{@{}p{0.66\linewidth}r@{}}
\hrow \wh{Input} & \wh{Value}\\
10-yr G-Sec par yield, w/e 20-Mar-26 & 6.85\% \\
\arow India sovereign default spread & 1.87\% \\
\trow Risk-free rate, default-free rupee & \textbf{4.98\%} \\
India equity risk premium & 7.08\% \\
\arow \quad of which mature market & 4.23\% \\
\quad of which India country risk & 2.85\% \\
\arow Unlevered beta, building materials & 0.900\\
\quad number of firms in the sample & 469\\
\arow Statutory tax rate, s.115BAA & 25.168\% \\
Long-run real GDP growth & 6.50\% \\
\arow Long-run inflation (RBI CPI target) & 4.00\% \\
\trow Nominal GDP, $(1+r)(1+i)-1$ & \textbf{10.76\%} \\
Terminal $g$ as \% of nominal GDP & 65.0\% \\
\trow Terminal growth rate & \textbf{6.994\%} \\
\end{tabular}
\end{extab}
\end{minipage}\hfill
\begin{minipage}[t]{0.48\textwidth}
\begin{extab}{Company-specific}{Harris--Pringle Case 2 applied throughout. Source:
company balance sheets and disclosed market capitalisations.}
\begin{tabular}{@{}p{0.50\linewidth}rr@{}}
\hrow \wh{} & \wh{Ambuja} & \wh{UltraTech}\\
Borrowings + leases (Rs cr) & 866 & 20,480\\
\arow Market capitalisation (Rs cr) & 99,095 & 316,633\\
Market debt / equity & 0.87\% & 6.47\% \\
\arow Debt / enterprise value & 0.87\% & 6.08\% \\
Levered beta $\beta_U(1+D/E)$ & 0.9079 & 0.9582\\
\arow Cost of equity, CAPM & 11.41\% & 11.76\% \\
Pre-tax cost of debt & 8.00\% & 8.34\% \\
\arow After-tax cost of debt & 5.99\% & 6.24\% \\
\trow \textbf{WACC} & \textbf{11.36\%} & \textbf{11.43\%} \\
\trow \textbf{Unlevered rate $\rho$} & \textbf{11.38\%} & \textbf{11.56\%} \\
\arow Regression beta, 60 monthly obs. & 1.151 & 1.086\\
$R^2$ of the regression & 0.267 & n/d\\
\arow WACC implied by market price & 8.60\% & 9.17\% \\
\end{tabular}
\end{extab}
\end{minipage}

\vspace{4mm}
\section*{E\quad Spread financial statements}

\begin{minipage}[t]{0.48\textwidth}
\begin{extab}{Ambuja: consolidated, as reported}{Source: FY2025-26 Integrated Annual
Report, consolidated statements. FY2024-25 is the company's restated figure.}
\begin{tabular}{@{}p{0.58\linewidth}rr@{}}
\hrow \wh{Rs crore} & \wh{FY25R} & \wh{FY26}\\
Revenue from operations & 33,989 & 40,446\\
\arow Government grants & 1,347 & 210\\
Revenue including grants & 35,336 & 40,656\\
\arow Other income & 2,654 & 834\\
Finance costs & 216 & 224\\
\arow Depreciation and amortisation & 2,297 & 3,570\\
PBT before exceptionals & 6,125 & 3,599\\
\arow Exceptional items & 21 & 301\\
Profit before tax & 6,104 & 3,299\\
\arow Profit, owners of parent & 4,303 & 4,728\\
Profit, non-controlling interest & 991 & 909\\
\trow \textbf{EBITDA, uniform definition} & \textbf{5,984} & \textbf{6,559}\\
\arow Property, plant and equipment & 25,049 & 33,801\\
Right-of-use assets & 1,465 & 1,483\\
\arow Capital work in progress & 9,793 & 9,085\\
Goodwill & 10,943 & 13,546\\
\arow Other intangible assets & 5,309 & 9,433\\
Equity, owners & 53,579 & 59,347\\
\arow Non-controlling interest & 10,368 & 12,499\\
Borrowings, non-current & 14 & 45\\
\arow Borrowings, current & 12 & 8\\
Lease liabilities & 762 & 813\\
\arow Cash and equivalents & 5,043 & 892\\
Other bank balances & 1,129 & 67\\
\arow Current investments & 1,822 & 0\\
Deferred tax liabilities, net & 2,433 & 3,466\\
\end{tabular}
\end{extab}
\end{minipage}\hfill
\begin{minipage}[t]{0.48\textwidth}
\begin{extab}{UltraTech: standalone ten-year highlights}{Source: FY2025-26 Integrated
and Sustainability Report, Standalone Financial Highlights. R = restated by the
company.}
\begin{tabular}{@{}p{0.34\linewidth}rrrrr@{}}
\hrow \wh{Rs crore} & \wh{FY17} & \wh{FY20} & \wh{FY23R} & \wh{FY25} & \wh{FY26}\\
Volume (Mn.T) & 48.9 & 77.5 & 100.1 & 125.8 & 140.8\\
\arow Revenue & 23,891 & 40,649 & 61,237 & 71,895 & 82,170\\
Operating expenses & 18,922 & 31,997 & 50,952 & 59,599 & 66,730\\
\arow EBITDA & 4,969 & 8,652 & 10,286 & 12,296 & 15,439\\
Other income & 660 & 727 & 505 & 693 & 376\\
\arow Depreciation & 1,268 & 2,455 & 2,773 & 3,739 & 4,055\\
EBIT & 4,361 & 6,924 & 8,018 & 9,250 & 11,761\\
\arow Interest & 571 & 1,704 & 756 & 1,465 & 1,630\\
Profit before tax & 3,790 & 5,220 & 7,262 & 7,785 & 10,131\\
\arow Tax expense & 1,148 & $-$236 & 2,310 & 1,504 & 2,622\\
Net profit & 2,628 & 5,456 & 4,951 & 6,193 & 7,405\\
\arow Net fixed assets & 24,387 & 49,486 & 63,077 & 83,764 & 89,161\\
Investments in subs & 746 & 5,990 & 3,187 & 12,999 & 12,541\\
\arow Liquid investments & 8,663 & 5,882 & 6,246 & 3,504 & 5,018\\
Net working capital & $-$841 & 190 & $-$3,140 & $-$2,031 & $-$3,071\\
\arow Net worth & 23,941 & 38,296 & 53,408 & 69,678 & 74,663\\
Loan funds & 6,240 & 18,282 & 8,750 & 19,460 & 19,617\\
\arow Lease liabilities & 0 & 893 & 953 & 901 & 863\\
Deferred tax liabilities & 2,774 & 4,077 & 6,258 & 8,198 & 8,506\\
\trow \textbf{Capital employed} & \textbf{32,955} & \textbf{61,548} & \textbf{69,369} & \textbf{98,236} & \textbf{103,649}\\
\end{tabular}
\end{extab}
\end{minipage}

\vspace{1mm}
{\scriptsize\color{ash}Capital employed ties both ways in every one of the ten years:
net worth plus loan funds plus leases plus deferred tax equals net fixed assets plus
investments plus liquid investments plus net working capital.\par}
\clearpage

% ---------------------------------------------------------------- F
\section*{F\quad Forecast driver tables}

\subsection*{F.1\quad Ambuja Cements: ten-year explicit forecast}
\vspace{1mm}
{\scriptsize\ambforecast\par}
\vspace{1mm}
{\scriptsize\color{ash}Consolidated. Volume is capacity times utilisation. Revenue is
volume times realisation. EBITDA is the sum of a standalone block and a subsidiary
block, each with its own margin. NOPAT is EBIT at the statutory 25.168\%. Invested
capital includes goodwill of Rs 22,979 crore held flat. Source: authors' model.\par}

\vspace{5mm}
\subsection*{F.2\quad UltraTech Cement: five-year explicit forecast}
\vspace{1mm}
{\footnotesize\utclforecast\par}
\vspace{1mm}
{\scriptsize\color{ash}Standalone core plus a subsidiary block. Volume is capacity
times utilisation; EBITDA is volume times EBITDA per tonne. Invested capital is the
consolidated operating asset base plus net working capital, and excludes the Rs 7,909
crore of consolidated goodwill. Source: authors' model.\par}

\vspace{5mm}
\subsection*{F.3\quad Terminal year, both companies}
\vspace{1mm}
\begin{minipage}[t]{0.48\textwidth}
\begin{extab}{Ambuja terminal year}{Source: authors' model.}
\begin{tabular}{@{}p{0.62\linewidth}r@{}}
\hrow \wh{Line} & \wh{Value}\\
NOPAT in the terminal year (Rs cr) & 8,908\\
\arow Closing operating invested capital & 74,610\\
Reinvestment $= g \times$ operating IC & 5,218\\
\arow \textbf{Terminal FCFF} & \textbf{3,690}\\
Implied terminal ROIC & 11.94\% \\
\arow Implied reinvestment rate ($g/$ROIC) & 58.6\% \\
Terminal value, Gordon (Rs cr) & 84,458\\
\arow Implied exit EV/EBITDA & 4.65x\\
Terminal ROIC less WACC & $+$0.58 pt\\
\end{tabular}
\end{extab}
\end{minipage}\hfill
\begin{minipage}[t]{0.48\textwidth}
\begin{extab}{UltraTech terminal year}{Source: authors' model.}
\begin{tabular}{@{}p{0.62\linewidth}r@{}}
\hrow \wh{Line} & \wh{Value}\\
NOPAT in the terminal year (Rs cr) & 17,920\\
\arow Closing invested capital & 112,483\\
Reinvestment $= g \times$ IC & 7,867\\
\arow \textbf{Terminal FCFF} & \textbf{10,053}\\
Implied terminal ROIC & 15.93\% \\
\arow Implied reinvestment rate ($g/$ROIC) & 43.9\% \\
Terminal value, Gordon (Rs cr) & 231,777\\
\arow Implied exit EV/EBITDA & 7.52x\\
Terminal ROIC less WACC & $+$4.50 pt \FAIL\\
\end{tabular}
\end{extab}
\end{minipage}

\vspace{4mm}
\subsection*{F.4\quad Sensitivity grids: value per share (Rs), WACC against terminal growth}
\vspace{1mm}
\begin{minipage}[t]{0.44\textwidth}
{\footnotesize\bfseries\color{navy}Ambuja Cements}\par\vspace{2pt}
{\footnotesize\ambgrid\par}
\vspace{2pt}{\scriptsize\color{ash}Traded price Rs 400.90. No cell reaches it.\par}
\end{minipage}\hfill
\begin{minipage}[t]{0.52\textwidth}
{\footnotesize\bfseries\color{navy}UltraTech Cement (rebuilt --- see Appendix B, item 1)}\par\vspace{2pt}
{\footnotesize\utclgrid\par}
\vspace{2pt}{\scriptsize\color{ash}Traded price Rs 10,745. No cell reaches it. The
11.43\% row is the base-case WACC.\par}
\end{minipage}
\clearpage

% ---------------------------------------------------------------- G, H
\section*{G\quad Method equivalence: why all five agree}

\begin{paracol}{2}
The five methods are one model rearranged, and they agree only if every consistency
holds. The mechanics are worth setting out, because the agreement is the proof that
the model is internally sound.

\textbf{FCFF at WACC} is the primary build: NOPAT plus depreciation less capex less
the increase in net working capital, discounted at the weighted average cost of
capital, plus a Gordon terminal value.

\textbf{CCF at $\rho$} adds the interest tax shield back into the cash flow and
discounts the whole capital cash flow at the unlevered rate. Because debt is set to
$D/V$ times enterprise value at every date, the shield grows with the firm and
carries the same risk as the assets, which is why $\rho$ is the correct rate.

\textbf{APV at $\rho$} separates the two: unlevered enterprise value plus the present
value of the tax shield. For Ambuja the shield is worth Rs 175 crore against an
enterprise value of Rs 39,833 crore — which is itself a finding. Ambuja creates
essentially no value through leverage, so any argument for the shares must rest
entirely on operations.

\textbf{FCFE at $K_e$} subtracts after-tax interest and adds net new borrowing.
Because a constant debt ratio is assumed, net new borrowing equals $g$ times debt in
the steady state, and the equity flows grow at the same rate as the firm. This is
where the most common error in the course appears — discounting a growing FCFE at a
constant $K_e$ while implicitly holding debt fixed — and the visible net borrowing
line is what prevents it.

\textbf{EVA at WACC} is opening invested capital plus the present value of future
economic value added. It reconciles only because FCFF was defined as NOPAT less the
change in invested capital, and because the terminal year is a true steady state.
Both conditions are necessary; drop either and the identity breaks.

\switchcolumn
\begin{extab}{Equivalence, both companies}{Source: authors' models, independently
reproduced.}
\begin{tabular}{@{}p{0.26\linewidth}p{0.16\linewidth}rr@{}}
\hrow \wh{Method} & \wh{Rate} & \wh{AMB} & \wh{UTCL}\\
FCFF & WACC & 99.42 & 5,024\\
\arow CCF & $\rho$ & 99.42 & 5,024\\
APV & $\rho$ & 99.42 & 5,024\\
\arow FCFE & $K_e$ & 99.42 & 5,024\\
EVA & WACC & 99.42 & 5,024\\
\trow \textbf{Max difference} & & \textbf{0.00} & \textbf{0.00}\\
\end{tabular}
\end{extab}

\begin{extab}{APV decomposition}{The tax shield is trivial for Ambuja and small for
UltraTech. Source: authors' models.}
\begin{tabular}{@{}p{0.50\linewidth}rr@{}}
\hrow \wh{Rs crore} & \wh{AMB} & \wh{UTCL}\\
Unlevered enterprise value & 39,658 & 163,214\\
\arow PV of interest tax shield & 175 & 4,554\\
\trow \textbf{Enterprise value} & \textbf{39,833} & \textbf{167,767}\\
Shield as \% of EV & 0.4\% & 2.7\% \\
\end{tabular}
\end{extab}

\begin{okbox}\boxtitleg{What we did to check this}
Both models were rebuilt from their input assumptions in a separate Python
implementation that shares no code with the workbooks. The rebuild reproduced
Ambuja's value to the paisa, UltraTech's to within 2.5 paise, Ambuja's published
sensitivity grid cell for cell, and the EVA-to-enterprise-value identity for both. A
model that only agrees with itself has not been checked.
\end{okbox}
\end{paracol}

\vspace{3mm}
\section*{H\quad Peer set and relative valuation data}

\begin{minipage}[t]{0.56\textwidth}
\begin{extab}{Peer set, TTM to June 2026 (Ambuja FY2025-26 consolidated)}{Every figure
from the company's own annual report and disclosed market capitalisation. Source:
company reports.}
\begin{tabular}{@{}p{0.30\linewidth}rrrrr@{}}
\hrow \wh{Rs crore} & \wh{Ambuja} & \wh{UTCL} & \wh{Shree} & \wh{JK Cem} & \wh{Ramco}\\
Market cap & 99,095 & 316,633 & 83,058 & 39,252 & 21,939\\
\arow Net debt & $-$93 & 16,444 & $-$6,157 & 5,717 & 3,644\\
Non-controlling int. & 12,499 & 0 & 0 & 0 & 0\\
\arow Enterprise value & 111,501 & 333,077 & 76,901 & 44,969 & 25,583\\
EBITDA & 6,559 & 15,817 & 4,037 & 2,284 & 1,348\\
\arow Revenue & 40,446 & 86,070 & 19,985 & 13,621 & 9,211\\
Normalised PAT & 2,260 & 7,762 & 1,534 & 1,096 & 178\\
\arow Net worth & 59,347 & 74,663 & 22,512 & 6,961 & 8,142\\
Shares (crore) & 247.18 & 29.47 & 3.61 & 7.73 & 23.63\\
\trow \textbf{EV/EBITDA} & \textbf{17.0} & \textbf{21.1} & \textbf{19.1} & \textbf{19.7} & \textbf{19.0}\\
\arow EV/Sales & 2.76 & 3.87 & 3.85 & 3.30 & 2.78\\
P/E, normalised & 43.9 & 40.8 & 54.2 & 35.8 & 122.9\\
\arow P/B & 1.67 & 4.24 & 3.69 & 5.64 & 2.69\\
ROE & 3.8\% & 10.4\% & 6.8\% & 15.7\% & 2.2\% \\
\arow Basis of accounts & Consol. & S/A & S/A & S/A & S/A\\
\end{tabular}
\end{extab}
\end{minipage}\hfill
\begin{minipage}[t]{0.40\textwidth}
\begin{extab}{Reverse cross-check: what an exit multiple implies}{UltraTech, WACC
11.43\%, nominal GDP 10.76\%. Source: authors' model.}
\begin{tabular}{@{}p{0.42\linewidth}rp{0.28\linewidth}@{}}
\hrow \wh{Exit EV/EB} & \wh{Implied $g$} & \wh{Verdict}\\
19.05x (peer median) & 9.8\% & 91\% of GDP\\
\arow 15.0x & 9.4\% & Implausible\\
12.0x & 8.8\% & Implausible\\
\arow 10.0x & 8.1\% & Stretching\\
8.0x & 6.8\% & Arguable\\
\arow \textbf{7.52x (our model)} & \textbf{6.99\%} & \textbf{Base case}\\
\end{tabular}
\end{extab}

\begin{warnbox}\boxtitler{Why the peer median is rejected rather than adopted}
A peer median EV/EBITDA of 19.4x, solved back through the steady-state identity,
requires perpetual growth of 10.99\% for Ambuja and 9.8\% for UltraTech, against
nominal GDP of 10.76\%. A company that grows faster than the economy forever
eventually becomes the economy. Adopting the peer multiple would require us to
contradict our own central argument, so we reject it and say why.
\end{warnbox}
\end{minipage}
\clearpage

% ---------------------------------------------------------------- I
\section*{I\quad Source log}

Every input cell in both models traces to a row below. Rows marked OPEN are items that
must be re-sourced before submission. No data anywhere in either workbook comes from
moneycontrol, screener.in, Yahoo Finance, investing.com, trendlyne, tijori or
Wikipedia.

\vspace{2mm}
{\footnotesize
\begin{tabular}{@{}p{0.028\textwidth}p{0.055\textwidth}p{0.215\textwidth}p{0.135\textwidth}p{0.37\textwidth}p{0.075\textwidth}@{}}
\hrow \wh{\#} & \wh{Co.} & \wh{Item} & \wh{Publisher} & \wh{Location} & \wh{Accessed}\\
1 & AMB & Consolidated P\&L and balance sheet & Ambuja Cements Ltd & Integrated Annual Report 2025-26, consolidated statements & 26-Aug-26\\
\arow 2 & AMB & Volume 74 Mn.T, revenue Rs 40,656 cr & Ambuja Cements Ltd & Same report, Growth Trajectory highlights & 26-Aug-26\\
3 & AMB & Capacity 109 MTPA, target 119 MTPA & Ambuja Cements Ltd & Same report, capacity disclosure & 26-Aug-26\\
\arow 4 & AMB & Market capitalisation Rs 99,095 cr & Ambuja Cements Ltd & Same report, highlights, as at 31-Mar-2026 & 26-Aug-26\\
5 & AMB & Holding in ACC 50.05\% & Ambuja Cements Ltd & Same report, subsidiary listing & 26-Aug-26\\
\arow 6 & AMB & Tax note: Rs 1,179.71 cr provisions released & Ambuja Cements Ltd & Same report, note 32(ii) & 26-Aug-26\\
7 & AMB & Duty and taxes paid under protest Rs 614.10 cr & Ambuja Cements Ltd & Same report, note 53, non-current & 26-Aug-26\\
\arow 8 & AMB & \textbf{OPEN:} merger swap ratios and 13.7\% dilution & Ambuja filings & Reached the model through broker research. Must be re-sourced from the scheme filing. & 26-Aug-26\\
9 & UTCL & Ten-year standalone financial highlights & UltraTech Cement Ltd & Integrated and Sustainability Report 2025-26, Standalone Financial Highlights & 26-Aug-26\\
\arow 10 & UTCL & Standalone and consolidated statements & UltraTech Cement Ltd & Same report, pp. 227--233 & 26-Aug-26\\
11 & UTCL & Capital work in progress Rs 7,871.47 cr & UltraTech Cement Ltd & Same report, note 02 & 26-Aug-26\\
\arow 12 & UTCL & India Cements: 23,24,16,830 shares, Rs 8,970.17 cr, 74.99\% & UltraTech Cement Ltd & Same report, non-current investments note & 26-Aug-26\\
13 & UTCL & Market cap Rs 316,633 cr; March 2026 NSE close Rs 10,745 & UltraTech Cement Ltd & Same report, Shareholder Information & 26-Aug-26\\
\arow 14 & UTCL & Q1 FY2026-27 unaudited standalone results & UltraTech Cement Ltd & q1fy27results.pdf, standalone results & 26-Aug-26\\
15 & UTCL & \textbf{OPEN:} capacity, utilisation, capex guidance, per-tonne split & UltraTech earnings call & Reached the model through a broker note relaying the call. Must be re-sourced from the transcript. & 26-Aug-26\\
\arow 16 & UTCL & \textbf{OPEN:} standalone capacity 172.5 Mn.T & Own derivation & Consolidated 192.3 less 5.4 overseas less India Cements 14.45. Not a disclosed figure. & 26-Aug-26\\
17 & Both & 10-year G-Sec par yield 6.85\%, w/e 20-Mar-2026 & RBI / FBIL & Ratios and Rates release dated 27-Mar-2026 & 27-Aug-26\\
\arow 18 & Both & India ERP 7.08\%; sovereign default spread 1.87\% & Damodaran, NYU Stern & Country Default Spreads and Risk Premiums, January 2026 & 26-Aug-26\\
19 & Both & Unlevered beta 0.90, building materials, cash-corrected & Damodaran, NYU Stern & Betas by Sector (Global), January 2026, 469 firms & 26-Aug-26\\
\arow 20 & Both & Statutory tax rate 25.168\% & Own assumption & Section 115BAA: 22\% plus 10\% surcharge plus 4\% cess & 26-Aug-26\\
21 & Both & Peer financials and market capitalisations & Peer annual reports & Shree Cement, JK Cement, Ramco Cements, own FY2025-26 reports & 26-Aug-26\\
\arow 22 & Both & Price series for the beta regression & Vendor series & Supplied with an adjusted-close column, a vendor convention rather than exchange bhavcopy format. Labelled vendor-sourced. & 26-Aug-26\\
\end{tabular}\par}

\vspace{3mm}
{\footnotesize\bfseries\color{navy}Authoritative landing pages}\par\vspace{1mm}
{\scriptsize
Ambuja financials: \texttt{ambujacement.com/investors/financials} \textbullet\
UltraTech: \texttt{ultratechcement.com/investors/investor-reports} \textbullet\
Exchange filings: \texttt{nseindia.com/companies-listing/corporate-filings-financial-results} \textbullet\
G-Sec yields: \texttt{rbi.org.in/Scripts/BS\_ViewBulletin.aspx} \textbullet\
ERP and default spreads: \texttt{pages.stern.nyu.edu/\textasciitilde adamodar/New\_Home\_Page/datacurrent.html} \textbullet\
Sector betas: \texttt{pages.stern.nyu.edu/\textasciitilde adamodar/New\_Home\_Page/datafile/Betas.html}\par
\vspace{1mm}
Deep links to individual PDFs are deliberately not given: they are unstable, and
citing a precise document URL that was not actually followed would be worse than
citing the page that was. Each report is identified by name, year and note reference
in the table above, which is sufficient to locate the exact figure.\par}
\clearpage

% ---------------------------------------------------------------- J, K
\begin{paracol}{2}
\section*{J\quad Method note and glossary}

\subsection*{J.1\quad The framework applied}

This report follows the theoretically consistent valuation framework taught in the
XLRI Business Analysis and Valuation course. Its central claim is that a valuation is
correct when it preserves five consistencies — assumptions, risk, asset, liability and
multiplier — and that when all five hold, every discounted cash flow variant returns
the same equity value. That equivalence is not a stylistic flourish; it is the test.
Where methods disagree, one of the five consistencies has been broken and the model is
wrong.

Three further rules shape what follows. Relative valuation is presented after the
discounted cash flow and never anchors it, because a multiple is the output of a
valuation rather than its cause. Every number traces to a company annual report, an
exchange filing, or a regulator or government source; no figure comes from a financial
data aggregator. And where intrinsic value diverges from market price, the report names
the catalysts that would close the gap rather than asserting that the market is wrong.

\subsection*{J.2\quad Glossary}

\textbf{FCFF} --- free cash flow to the firm: NOPAT plus depreciation less capital
expenditure less the increase in net working capital. The cash available to all
providers of capital.

\textbf{NOPAT} --- net operating profit after tax: EBIT taxed at the marginal statutory
rate, with no deduction for interest, so it measures operating performance independent
of financing.

\textbf{WACC} --- the weighted average cost of capital, at market-value weights.

\textbf{$\rho$} --- the unlevered cost of capital: the return the assets would require
if the firm carried no debt. Under constant rebalancing, $\rho = \text{WACC} + K_d
t(D/V)$.

\textbf{CCF} --- capital cash flow: FCFF plus the interest tax shield, discounted at
$\rho$.

\textbf{APV} --- adjusted present value: unlevered value plus the present value of the
tax shield, valued separately.

\textbf{FCFE} --- free cash flow to equity: FCFF less after-tax interest plus net new
borrowing, discounted at the cost of equity.

\textbf{EVA} --- economic value added: NOPAT less a capital charge of WACC on opening
invested capital. Enterprise value equals opening invested capital plus the present
value of all future EVA.

\textbf{ROIC} --- return on invested capital: NOPAT over invested capital. Reported
both including and excluding goodwill, because the two answer different questions.

\textbf{Harris--Pringle (Case 2)} --- the debt policy in which the firm maintains a
constant debt-to-value ratio through continuous rebalancing. Under it the tax shield
carries asset risk and is discounted at $\rho$, and beta is relevered as
$\beta_U(1+D/E)$ with no tax term.

\switchcolumn
\section*{K\quad Limitations and disclosures}

\subsection*{K.1\quad Limitations we accept}

\subsubsection{Ambuja's history is two years deep} The consolidated return trend is
asserted from two points, one of which is restated. This is the single largest
weakness in Part I.

\subsubsection{One open sourcing item each} Ambuja's merger terms and UltraTech's
capacity, utilisation and capex guidance reached the models through broker research,
which our policy does not admit. Both are flagged rather than quietly used, and
UltraTech's standalone capacity of 172.5 Mn.T is a derivation rather than a
disclosure.

\subsubsection{The two models are not on a common basis} One is consolidated and one
is a standalone core with subsidiaries added. The reason is data availability, the
effect is quantified in Appendix A, and the fix is to consolidate both.

\subsubsection{The beta regression is weak} Ambuja's regression explains 27\% of
variance, and the price series carries a vendor adjusted-close convention rather than
exchange bhavcopy format. Bottom-up betas are used for both companies, and in both
cases that choice raises the valuation rather than lowering it.

\subsubsection{Both carry a multiplier PARTIAL} The implied terminal multiples sit far
below anything cement has traded at. We have taken a documented position on that gap
rather than resolved it, and a reader who takes the opposite view should say so
explicitly rather than average the two answers.

\subsubsection{UltraTech's terminal return does not fade} Recorded as a FAIL on Golden
Check 6. Correcting it would lower the valuation by roughly Rs 464 per share, so the
error runs against our own recommendation.

\subsection*{K.2\quad On the use of AI assistance}

AI tools were used for data gathering, for spreading and cross-checking the financial
statements, for the independent Python re-implementation used to audit both workbooks,
and for drafting. Every figure in this document remains human-verifiable through the
source log in Appendix I and reproducible from the two accompanying Excel models. The
assumptions are the authors' own judgements, and the errors found in Appendix B were
identified during the build rather than reported by a reviewer afterwards. No source,
transcript quote or figure has been fabricated; where an input could not be obtained
from an authorised source it is marked as an assumption or as an open item.

\subsection*{K.3\quad Basis and disclaimer}

Valuation date 31 March 2026, with market data at the 30 March close — 31 March was
not a trading day. All figures are in rupees crore unless stated. Ambuja is presented
on a consolidated basis; UltraTech on a standalone core with subsidiary holdings added
at the bridge. Financial year references are Indian fiscal years ending 31 March.

This document was prepared for an academic valuation exercise. It is not investment
research within the meaning of any securities regulation, not a recommendation to buy
or sell any security, and not advice. The authors hold no position in any company
mentioned.

\begin{keybox}\boxtitle{Reproducibility}
Every exhibit in Parts I, II and III is generated directly from the accompanying
models. The two corrected Excel workbooks, the chart-generation scripts and the
independent Python re-implementation used for the audit in Appendix B are delivered
alongside this document.
\end{keybox}
\end{paracol}
